% =====================================================================
% Слайд 13 — Математика RL-слоя: формулы (state, action, reward)
% =====================================================================
\begin{frame}{Математика RL-слоя: состояние, действие, награда}
  \footnotesize\setlength{\abovedisplayskip}{3pt}
  \setlength{\belowdisplayskip}{3pt}

  \textbf{Состояние} $s_t \in \mathbb{R}^{18}$ (\rlhi{не зависит от
  длины пьесы}):
  \[
    s_t = \bigl(\,\hat\alpha_t,\ \tau_{t-K:t},\ e_{t-K:t},\
                 \hat\varphi(t)/N\,\bigr)
  \]
  \begin{itemize}\setlength{\itemsep}{0.15em}
    \item $\hat\alpha_t$ --- 7-числовая сжатка HSMM-распределения
          (max, entropy, argmax/$N$, top-3 mass, 3 рел.\ индекса);
    \item $\tau_{t-K:t}$ --- последние $K$ оценок темпа,\quad
          $e_{t-K:t}$ --- эмиссионные ошибки;
    \item $\hat\varphi(t)/N$ --- относительная позиция в партитуре.
  \end{itemize}

  \textbf{Действие:} $a_t \in [0.5,\, 2.0]$ --- темповый коэффициент,
  отдаётся каждые 50\,мс.

  \vspace{0.3em}
  \textbf{Награда} (формула\,(1) тезисов):
  \[
    r_t \;=\;
      \underbrace{-\bigl|\,t^{\text{render}}_t - t^{\text{perf}}_t\bigr|}_{\text{рассинхр.}}
      \;
      \underbrace{-\,\lambda\,L_{\text{align}}(t)}_{\text{ошибка трекера}}
      \;
      \underbrace{-\,\mu\,(a_t - a_{t-1})^{2}}_{\text{резкий перепад}}
  \]
  \begin{itemize}\setlength{\itemsep}{0.15em}
    \item 1-й член --- основной сигнал рассинхронизации рендера и исполнителя;
    \item 2-й --- мягкая регуляризация: не уезжать от базового трекера;
    \item 3-й --- гладкость темпа, чтобы оркестр не «дёргался».
  \end{itemize}
  $\lambda, \mu$ подбираются на PPO по ASAP.
\end{frame}

% =====================================================================
% Слайд 13b — Энкодер состояния s_t (TikZ-схема)
% =====================================================================
\begin{frame}{Энкодер состояния $s_t$}
  \small

  Все четыре источника сжимаются в один фиксированный вектор
  $s_t \in \mathbb{R}^{18}$ --- который подаётся в MLP $2\!\times\!64$.

  \vspace{0.6em}
  \begin{center}
    \resizebox{0.92\textwidth}{!}{%
      % СХЕМА 9 — Энкодер состояния s_t
\begin{tikzpicture}[
  every node/.style={font=\scriptsize, align=center},
  src/.style={draw, thick, minimum height=0.65cm, minimum width=2.0cm,
              fill=white, inner sep=2pt},
  feat/.style={draw, thick, minimum height=0.55cm, minimum width=1.8cm,
               fill=black!4, inner sep=2pt},
  state/.style={draw, very thick, minimum height=0.9cm, minimum width=2.2cm,
                fill=white, inner sep=2pt},
  arr/.style={-{Latex[length=2mm]}, thick},
]
  \node[src] (a) at (0, 2.0) {HSMM $\alpha_t$};
  \node[src] (b) at (0, 1.0) {Tempo Tracker};
  \node[src] (c) at (0, 0.0) {Emission error};
  \node[src] (d) at (0, -1.0) {position/$N$};

  \node[feat] (af) at (3.0, 2.0) {AlphaSummary (7)};
  \node[feat] (bf) at (3.0, 1.0) {$\tau_{t-K:t}$ (5)};
  \node[feat] (cf) at (3.0, 0.0) {$e_{t-K:t}$ (5)};
  \node[feat] (df) at (3.0, -1.0) {$\hat\varphi/N$ (1)};

  \node[state] (s) at (6.5, 0.5) {$s_t \in \mathbb{R}^{18}$};

  \draw[arr] (a) -- (af);
  \draw[arr] (b) -- (bf);
  \draw[arr] (c) -- (cf);
  \draw[arr] (d) -- (df);
  \foreach \f in {af, bf, cf, df}
    \draw[arr] (\f.east) -- (s.west);

  \draw[arr] (s.east) -- ++ (1.3, 0)
       node[right, font=\scriptsize] {$\pi_\theta$ (MLP $2{\times}64$)};
\end{tikzpicture}
%
    }
  \end{center}

  \vspace{0.4em}
  \begin{columns}[T,onlytextwidth]
    \column{0.50\textwidth}
    \footnotesize
    \textbf{Что важно:}
    \begin{itemize}\setlength{\itemsep}{0.2em}
      \item размер $s_t$ \rlhi{не зависит} от длины пьесы;
      \item $\hat\alpha_t$ --- сжатка HSMM, а не сам $\alpha$
            (иначе $\dim s_t \to N$);
      \item $K{=}5$ --- баланс между контекстом и инерцией.
    \end{itemize}

    \column{0.50\textwidth}
    \footnotesize
    \textbf{Что это даёт:}
    \begin{itemize}\setlength{\itemsep}{0.2em}
      \item одна и та же политика для пьес 200\ldots4000 state;
      \item MLP инициализируется в BC \emph{один раз} и переносится
            между корпусами;
      \item признаки физически интерпретируемые.
    \end{itemize}
  \end{columns}
\end{frame}
